\section{Task 1.4 — Explainability (20\,pts)}
We analyse the trained CNN with three complementary methods on validation images
(Fig.~\ref{fig:explain}): vanilla-gradient saliency, Grad-CAM on the last
convolutional block, and occlusion sensitivity (sliding a neutral patch and
measuring the drop in P(AI)). We examine the most borderline true-positive (AI),
true-negative (real), false-positive and false-negative cases.

\textbf{Findings.} Saliency is noisy and spatially \emph{distributed} rather than
tied to one object, and Grad-CAM frequently highlights broad background/sky
regions instead of the main subject — consistent with a detector that keys on
global texture statistics rather than semantics. The occlusion analysis is the
most revealing: replacing a patch with neutral grey almost always \emph{raises}
P(AI) (negative $\Delta$P(AI) in Fig.~\ref{fig:explain}), i.e.\ the model reads
flat, low-detail regions as AI-like. Averaged over many images
(Fig.~\ref{fig:occ}) this effect is strongest in the image centre (mean
$\Delta$P(AI) centre vs.\ border $\OccCenterReal$ vs.\ $\OccBorderReal$ for real,
$\OccCenterAi$ vs.\ $\OccBorderAi$ for AI), exposing two biases worth flagging: a
\emph{smoothness confound} (AI images are smoother, so the model equates
smoothness with ``fake'') and a mild \emph{centre bias}. These also explain the
errors: \emph{false positives} are smooth, low-detail real photos, and
\emph{false negatives} are unusually detailed AI images. We therefore treat the
explanations critically — the model relies on a genuine but shortcut-prone cue
(local detail/smoothness) rather than a single robust artefact, which is exactly
why the Task~1.3 augmentation (which perturbs detail and compression) is needed.

\begin{figure}[t]
  \centering
  \includegraphics[width=0.86\linewidth]{explain_examples.png}
  \caption{Per-case explanations (rows: borderline TP-AI, TN-real, FP, FN).
  Columns: input, saliency, Grad-CAM, occlusion $\Delta$P(AI).}
  \label{fig:explain}
\end{figure}
\begin{figure}[t]
  \centering
  \begin{subfigure}{0.4\linewidth}
    \includegraphics[width=\linewidth]{baseline_importance.png}
    \caption{Classical feature importance.}\label{fig:imp}
  \end{subfigure}\hfill
  \begin{subfigure}{0.52\linewidth}
    \includegraphics[width=\linewidth]{explain_occlusion_avg.png}
    \caption{Mean occlusion importance (real vs.\ AI).}\label{fig:occ}
  \end{subfigure}
  \caption{What the models rely on: high-frequency spectral bins (a) and a mild
  centre bias in spatial attention (b).}
\end{figure}
